\chapter*{译者说明}
\addcontentsline{toc}{chapter}{译者说明}

\textbf{底本。}本书是 Dragan Miličić 的讲义《Localization and Representation Theory of Reductive Lie Groups》，未正式出版，作者自排（pdf\TeX{} 1.40.18，2017 年 10 月 22 日生成），印刷页 129 页，含 6 章与参考文献。底本 PDF 共 133 页，其中前 4 页为封面、空白页与目录，故\textbf{印刷页 = PDF 页 $-4$}。译文中的页码标注一律用印刷页，并在资料卡中同时给出 PDF 页。

\textbf{体例。}译本保留原书的章节结构与定理编号（「节.序号」）。原书定理、引理、命题、推论共用一个计数器，译本沿用这一约定。译者所加的解释性小字块以「校注」标出，其内容或回引作者前作《Lectures on algebraic D-modules》（本库编号 SY07-10，已有中文校订批注本），或补充原书略去的验证，或标出原书笔误；每条校注都注明依据。

\textbf{术语。}术语依《工程文件/work/术语表.tsv》统一，正文、书末术语索引与该表三处同步。有通行中译的用中译，无通行中译的保留英文原名（如 Verma 模、Kazhdan--Lusztig 多项式、Harish-Chandra 模）。人名一律保留拉丁字母。

\textbf{省略号。}底本中作者常以「\dots」省去具体的引文编号（如「由 \dots 推出」），译文照录为「\dots\dots」，不代补编号；凡能由上下文唯一确定的，另以校注说明。

\textbf{记号。}译本的记号尽量与底本一致；底本未定义而由译者补足的记号，在首次出现处用校注说明。
